\documentclass[11pt]{letter}
\usepackage[margin=1in]{geometry}

% Replace these macros with author-confirmed facts before submission.
\newcommand{\AuthorName}{FULL AUTHOR NAME}
\newcommand{\AuthorAffiliation}{FULL AFFILIATION}
\newcommand{\AuthorEmail}{EMAIL ADDRESS}
\newcommand{\AuthorPostalAddress}{POSTAL ADDRESS}

\signature{\AuthorName\\\AuthorAffiliation\\\AuthorEmail}
\address{\AuthorName\\\AuthorAffiliation\\\AuthorPostalAddress\\\AuthorEmail}

\begin{document}
\begin{letter}{Editors\\International Journal of Industrial Organization}
\opening{Dear Editors,}

Please consider the manuscript ``When Supplier Competition Worsens Policy Coordination: Procurement Rents and Regional Industrial-Policy Specialization'' for publication in the \emph{International Journal of Industrial Organization}.

The paper studies a two-layer industrial-organization problem. Two regional governments choose whether to take a locally attractive lead/buyer-side role or a complementary supplier-side role. When the regions specialize, suppliers compete in a procurement market. The procurement comparative static itself is standard: stronger supplier competition lowers expected minimum cost and, under a decreasing-reversed-hazard-rate condition, reduces expected winning supplier rent. The paper's contribution begins at the policy layer. Competition raises the total value of complementary specialization while reducing the locally captured rent that makes a region willing to take the supplier-side role. The coordinated differentiation threshold therefore falls with supplier competition while the decentralized threshold rises. Sufficiently strong competition can eliminate differentiated Nash equilibria even when differentiation remains the coordinated optimum.

We believe the manuscript is a natural fit for IJIO because its central comparative static is a market-structure primitive---the number of suppliers---and the result concerns how supplier-market competition changes the equilibrium of industrial policy.

A related companion manuscript, ``Industrial Policy Composition and Regional Value Chains,'' is disclosed transparently. The companion studies a continuous fixed-capacity regional policy-composition game and takes complementary surplus and local incidence largely as primitives. The present paper instead studies binary regional roles with an explicit private-information procurement continuation. At a fixed supplier count, the switching logic maps into the companion's reduced-form argument; this paper does not claim that fixed-market wedge as new. The new result is the supplier-market comparative static that moves coordinated and decentralized thresholds in opposite directions and can eliminate differentiated equilibria.

The manuscript is theoretical and uses no empirical data. Reproducibility materials contain Python/SymPy algebraic regression checks and figure-generation code. The manuscript contains a declaration of generative-AI and AI-assisted technologies and separately describes AI-assisted verification code.

\textbf{Administrative completion before submission:} Insert the author-confirmed statements on exclusivity, preprint status, funding, and competing interests here if the live portal or editor requires them in the cover letter.

Thank you for your consideration.

\closing{Sincerely,}
\end{letter}
\end{document}
